\documentclass[12pt,a4paper]{report}

% --- Cấu hình Tiếng Việt cho pdfLaTeX ---
\usepackage[utf8]{inputenc}
\usepackage[T5]{fontenc}
\usepackage[vietnamese]{babel}
\usepackage{times} 

% --- THIẾT LẬP LỀ TRANG ---
\usepackage{geometry}
\geometry{a4paper, top=2.5cm, bottom=2cm, left=2.5cm, right=2cm}

% --- THIẾT LẬP CỠ CHỮ CƠ BẢN 13PT ---
\usepackage[fontsize=13pt]{scrextend}

\usepackage[x11names,dvipsnames,svgnames,table]{xcolor} 
\usepackage{graphicx}
\usepackage[space]{grffile} 
\usepackage{float} 
\usepackage{amsmath} 
\usepackage[hidelinks]{hyperref}
\usepackage{booktabs} 
\usepackage{longtable} 
\usepackage{array}

% --- CẤU HÌNH HEADERS ---
\usepackage{titlesec}
\titleformat{\chapter}[display]
{\centering\bfseries\fontsize{18pt}{22pt}\selectfont}
{\chaptertitlename\ \thechapter}
{10pt}
{\fontsize{16pt}{20pt}\selectfont}

\titleformat{\section}
{\bfseries\fontsize{14pt}{16.8pt}\selectfont}
{\thesection}
{1em}
{}

\titleformat{\subsection}
{\bfseries\fontsize{13pt}{15.6pt}\selectfont}
{\thesubsection}
{1em}
{}

% --- KHOẢNG CÁCH TIÊU ĐỀ ---
\titlespacing*{\chapter}{0pt}{-20pt}{20pt}
\titlespacing*{\section}{0pt}{12pt}{6pt}
\titlespacing*{\subsection}{0pt}{12pt}{6pt}

% --- THỤT ĐẦU DÒNG VÀ KHOẢNG CÁCH ĐOẠN ---
\usepackage{indentfirst}
\setlength{\parindent}{0.5cm}
\setlength{\parskip}{12pt}

% --- CẤU HÌNH CAPTION ---
\usepackage{caption}
\DeclareCaptionFont{size11}{\fontsize{11pt}{13.2pt}\selectfont}
\captionsetup{font={it, size11}}

% --- TRANG BÌA ---
\usepackage{tikz}
\usepackage{setspace}

% --- CẤU HÌNH HEADER VÀ FOOTER ---
\usepackage{fancyhdr}
\pagestyle{fancy}
\fancyhf{}
\fancyhead[R]{\fontsize{8pt}{9.6pt}\selectfont Hồ Quốc Thắng 23129398}
\fancyfoot[C]{\fontsize{8pt}{9.6pt}\selectfont Bảo quản \& chế biến thuỷ sản | Aquatic Products Preservation and Processing Technologies | Nguyễn Anh Trinh}
\fancyfoot[R]{\fontsize{8pt}{9.6pt}\selectfont \thepage}
\renewcommand{\headrulewidth}{0.4pt}
\renewcommand{\footrulewidth}{0pt}

\graphicspath{{E:/study/KNVSTP/}}

\begin{document}
	
	% --- TRANG BÌA ---
	\pagenumbering{gobble}
	\pagestyle{empty}
	\begin{titlepage}
		\begin{tikzpicture}[remember picture, overlay]
			% Khung ngoài
			\draw[line width=2.5pt]
			([xshift=1.2cm, yshift=-1.2cm]current page.north west)
			rectangle
			([xshift=-1.2cm, yshift=1.2cm]current page.south east);
			% Khung trong
			\draw[line width=0.8pt]
			([xshift=1.45cm, yshift=-1.45cm]current page.north west)
			rectangle
			([xshift=-1.45cm, yshift=1.45cm]current page.south east);
		\end{tikzpicture}
		
		\begin{center}
			
			\vspace*{0.4cm}
			{\fontsize{13pt}{16pt}\selectfont\textbf{BỘ GIÁO DỤC VÀ ĐÀO TẠO}}\\
			{\fontsize{13pt}{16pt}\selectfont\textbf{TRƯỜNG ĐẠI HỌC NÔNG LÂM TP.HCM}}\\
			{\fontsize{13pt}{16pt}\selectfont\textbf{KHOA CÔNG NGHỆ HÓA HỌC VÀ THỰC PHẨM}}
			
			\vspace{0.6cm}
			
			\includegraphics[width=4.5cm]{0.jpg}
			
			\vspace{0.8cm}
			
			{\fontsize{18pt}{22pt}\selectfont\textbf{BÀI TẬP VỀ NHÀ}}
			
			\vspace{0.4cm}
			
			{\fontsize{16pt}{20pt}\selectfont\textbf{BẢO QUẢN \& CHẾ BIẾN THỦY SẢN}}
			
			\vspace{1.6cm}
			
		\end{center}
		
		\noindent
		{\fontsize{13pt}{16pt}\selectfont\textbf{GVGD:} Nguyễn Anh Trinh}\\
		{\fontsize{13pt}{16pt}\selectfont\textbf{}}
		
		\vspace{0.8cm}
		
		\begin{center}
			{\fontsize{13pt}{16pt}\selectfont\textbf{Sinh viên thực hiện}}\\[6pt]
			\begin{tabular}{@{\hspace{0.5cm}} p{6cm} @{\hspace{1.5cm}} p{3cm} @{\hspace{0.5cm}}}
				\toprule
				\textbf{Họ và tên} & \textbf{MSSV} \\
				\midrule
				Hồ Quốc Thắng     & 23129398 \\[3pt]
				
				\bottomrule
			\end{tabular}
		\end{center}
		
		\vfill
		
		\begin{center}
			{\fontsize{13pt}{16pt}\selectfont\textbf{Thành phố Hồ Chí Minh, tháng 3/2026}}
		\end{center}
		
		
	\end{titlepage}
	
	% --- TRANG MỤC LỤC ---
	\renewcommand{\contentsname}{MỤC LỤC}
	\tableofcontents
	\newpage

\pagenumbering{arabic}
\setcounter{page}{1}



\section{Mở đầu}

Sự gia tăng nhu cầu đối với các sản phẩm thủy hải sản chế biến sẵn đòi hỏi việc phát triển những quy trình bảo quản tối ưu nhằm duy trì chất lượng cảm quan và giá trị dinh dưỡng. Bạch tuộc (\textit{Octopus sp.}) là một đối tượng thủy sản có đặc tính sinh học và cấu trúc cơ thịt phức tạp, khác biệt hoàn toàn với các loài cá xương truyền thống. Với hàm lượng protein cao, ít béo nhưng giàu các axit béo chưa no đa bậc (PUFA), bạch tuộc cực kỳ nhạy cảm với các biến đổi tự phân và oxy hóa sau thu hoạch \cite{lipid_storage}. Việc xây dựng một kế hoạch nghiên cứu và thực thi bảo quản toàn diện từ khâu sơ chế đến đóng gói thành phẩm không chỉ đảm bảo an toàn thực phẩm mà còn tạo ra giá trị gia tăng đáng kể cho nguồn lợi hải sản này.

\section{Phân tích đặc tính sinh học và cấu trúc cơ thịt bạch tuộc}

Cơ thịt của các loài động vật thân mềm chân đầu (Cephalopods) nói chung và bạch tuộc nói riêng có cấu trúc sợi cơ đặc biệt, được bao bọc bởi một hệ thống mô liên kết dày đặc với các liên kết chéo collagen bền vững hơn so với cá \cite{octopus_texture}. Chính cấu trúc này tạo nên độ dai giòn đặc trưng, nhưng đồng thời cũng là thách thức trong việc làm mềm và thẩm thấu các chất bảo quản như muối. Sự phân bố của mô liên kết không đồng đều giữa các bộ phận giải phẫu như đầu, thân và râu dẫn đến tốc độ biến đổi lý hóa khác nhau trong suốt quá trình xử lý nhiệt hoặc sấy khô \cite{octopus_drying}.

Về mặt sinh hóa, bạch tuộc có hoạt tính tự phân (autolysis) rất cao. Sau khi chết, hệ thống enzyme nội tại bắt đầu phân giải protein và các hợp chất nitrogen phi protein, tạo điều kiện thuận lợi cho vi khuẩn phát triển \cite{octopus_oregano}. Lớp dịch nhớt (slime) bao phủ bên ngoài da bạch tuộc chứa nồng độ glycoprotein cao và là môi trường nuôi dưỡng vi sinh vật sơ cấp. Do đó, việc loại bỏ nhớt triệt để là bước đi tiên quyết để ngăn chặn sự hư hỏng ngay từ giai đoạn đầu \cite{squid_acids}.

\section{Quy trình sơ chế và xử lý tiền chế biến (Preliminary Processing)}

Giai đoạn sơ chế đóng vai trò thiết lập nền tảng cho sự thành công của toàn bộ quy trình bảo quản. Mục tiêu của giai đoạn này là loại bỏ các cơ quan dễ bị hư hỏng nhất và làm sạch bề mặt nguyên liệu để tối ưu hóa khả năng thẩm thấu muối ở bước tiếp theo.

\subsection{Tiếp nhận và xử lý cơ học ban đầu}

Bạch tuộc tươi được lựa chọn dựa trên các tiêu chí về độ săn chắc của thịt và màu sắc tự nhiên của lớp sắc tố da. Quá trình làm sạch và cắt bỏ nội tạng (gutting) phải đảm bảo tính triệt để. Các bộ phận như túi mực, hệ thống tiêu hóa, mắt và răng (beak) cần được loại bỏ hoàn toàn \cite{octopus_methods}. Việc làm vỡ túi mực không chỉ ảnh hưởng đến màu sắc mà còn có thể thúc đẩy nhanh quá trình oxy hóa lipid do sự hiện diện của các hợp chất chứa sắt trong mực \cite{lipid_storage}.

\subsection{Kỹ thuật xử lý dịch nhớt và cải thiện cấu trúc}

Bạch tuộc có lớp nhớt dày và dính, đóng vai trò như một rào cản vật lý ngăn cản sự thẩm thấu của muối và sự thoát hơi nước trong quá trình sấy \cite{squid_acids}. Hai phương pháp chính được áp dụng:

\begin{enumerate}
    \item \textbf{Sử dụng muối khô (Dry Salting):} Việc chà xát trực tiếp muối hạt lên bề mặt da bạch tuộc tạo ra một áp suất thẩm thấu cục bộ cực cao, làm đông tụ các glycoprotein trong nhớt và tách chúng ra khỏi bề mặt da.
    \item \textbf{Sử dụng dung dịch axit hữu cơ:} Ngâm bạch tuộc trong axit citric hoặc axit lactic nồng độ 1-2\% làm giảm pH bề mặt, ức chế vi sinh vật \cite{squid_acids}.
\end{enumerate}

\begin{table}[H]
\centering
\caption{Thông số sơ chế bạch tuộc đề xuất}
\begin{tabular}{lp{5cm}p{6cm}}
\toprule
\textbf{Thông số} & \textbf{Phương pháp đề xuất} & \textbf{Tác động} \\
\midrule
Loại bỏ nội tạng & Thủ công & Loại bỏ nguồn enzyme và vi khuẩn sơ cấp \cite{octopus_methods} \\
Tẩy nhớt & Muối khô 3\% + Tumbling & Làm sạch bề mặt, tăng thẩm thấu \cite{octopus_patent} \\
Làm mềm thịt & Axit hữu cơ 1-2\% & Giảm pH, tăng hòa tan collagen \cite{octopus_texture} \\
Rửa sạch & Nước tuần hoàn & Loại bỏ cặn bẩn trong giác bám \cite{octopus_patent} \\
\bottomrule
\end{tabular}
\end{table}

\section{Công nghệ muối bạch tuộc (Salting)}

\subsection{Lựa chọn phương pháp muối ướt (Wet Salting)}

Dựa trên đặc tính cơ thịt ít béo, phương pháp muối ướt (wet salting) được ưu tiên để duy trì sự đồng nhất về độ ẩm bề mặt, tránh hiện tượng da bị nhăn nheo quá mức ($shrunken$ $skin$) \cite{wet_salting}. Dung dịch nước muối bão hòa (20-25\%) cho phép kiểm soát tốt tốc độ thẩm thấu.

\subsection{Động học của hoạt độ nước ($a_w$) trong bảo quản}

Hoạt độ nước ($a_w$) là thông số quyết định sự ổn định. Bạch tuộc tươi có $a_w \approx 1.0$ \cite{aw_test}. Các mốc $a_w$ cần đạt:
\begin{itemize}
    \item \textbf{Ngưỡng $a_w < 0.86$:} Ngăn chặn vi khuẩn gây bệnh như \textit{Staphylococcus aureus} \cite{aw_test}.
    \item \textbf{Ngưỡng $a_w < 0.75$:} Ức chế nấm men và vi khuẩn ưa muối \cite{aw_test}.
    \item \textbf{Ngưỡng $a_w < 0.60$:} Ổn định tuyệt đối về vi sinh \cite{aw_test}.
\end{itemize}

\section{Công nghệ sấy bạch tuộc (Drying)}

Sấy là quá trình hạ thấp hàm lượng nước đến mức mà các phản ứng sinh hóa đình trệ. Khuếch tán nội ($internal$ $diffusion$) là yếu tố kiểm soát tốc độ sấy bạch tuộc \cite{octopus_drying}.

\subsection{Hiện tượng tạo màng (Case Hardening)}

Khi sấy ở nhiệt độ quá cao (trên $60^\circ C$) ngay từ đầu, tốc độ bay hơi bề mặt vượt quá tốc độ khuếch tán nội, tạo lớp vỏ khô cứng ($case$ $hardening$) ngăn cản sự thoát nước bên trong, tạo điều kiện cho vi khuẩn kỵ khí phát triển \cite{case_hardening}. Kiểm soát nhiệt độ ở mức $45-50^\circ C$ trong giai đoạn đầu là cực kỳ quan trọng \cite{vacuum_drying}.

\subsection{So sánh các phương pháp sấy}

\begin{itemize}
    \item \textbf{Sấy tự nhiên:} Chi phí thấp nhưng khó kiểm soát vệ sinh.
    \item \textbf{Sấy năng lượng mặt trời hỗn hợp (MSD):} Tăng tốc độ sấy lên 74\% và giảm biến nâu \cite{octopus_drying}.
    \item \textbf{Sấy đối lưu cưỡng bức:} Đưa $a_w$ xuống 0.63 trong vòng 10 giờ tại $70^\circ C$ \cite{octopus_kinetics}.
\end{itemize}

\begin{table}[H]
\centering
\caption{Thông số động học sấy bạch tuộc}
\begin{tabular}{lll}
\toprule
\textbf{Thông số} & \textbf{Giá trị} & \textbf{Ý nghĩa} \\
\midrule
Nhiệt độ sấy tối ưu & $45-55^\circ C$ & Tránh case hardening \cite{vacuum_drying} \\
Độ ẩm thành phẩm & $15-20\%$ & Ngăn chặn nấm mốc ($a_w \approx 0.6$) \cite{fish_preservation} \\
Hệ số khuếch tán ($D_{eff}$) & $10^{-9}$ $m^2/s$ & Đặc trưng khả năng thoát nước \cite{pressure_drying} \\
\bottomrule
\end{tabular}
\end{table}

\section{Công nghệ xông khói bạch tuộc (Smoking)}

\subsection{Cơ chế bảo quản và tạo màu}

Xông khói tác động qua ba cơ chế: chống oxy hóa (phenol), kháng khuẩn (axit hữu cơ) và tạo cảm quan (phản ứng Maillard giữa carbonyl và amino) \cite{smoke_antioxidant, hun_khoi_ky_thuat}.

\subsection{Lựa chọn nguồn gỗ xông khói}

Tại Việt Nam, có thể tận dụng gỗ Nhãn (\textit{Longan}) và Vải (\textit{Lychee}) có hàm lượng phenol cao và khả năng kháng oxy hóa ($IC_{50}$ thấp) ấn tượng \cite{lychee_pulp}. Gỗ Sồi (\textit{Oak}) là lựa chọn phổ biến nhất cho màu vàng sẫm đẹp mắt \cite{go_soi_xong_khoi}.

\section{Đánh giá chất lượng và tiêu chuẩn an toàn (TCVN 5649:2006)}

Sản phẩm phải tuân thủ các giới hạn về kim loại nặng và vi sinh vật theo tiêu chuẩn quốc gia.

\begin{table}[H]
\centering
\caption{Giới hạn vi sinh vật (TCVN 5649:2006)}
\begin{tabular}{lc}
\toprule
\textbf{Loại vi sinh vật} & \textbf{Mức tối đa (trong 1g)} \\
\midrule
Tổng số vi sinh vật hiếu khí & $10^6$ \\
Coliforms & $10^2$ \\
E. coli & 10 \\
Salmonella & 0 (trong 25g) \\
Staphylococcus aureus & $10^2$ \\
Tổng số bào tử nấm men - nấm mốc & $10^3$ \\
\bottomrule
\end{tabular}
\end{table}

\section{Kết luận và khuyến nghị}

Quy trình bảo quản bạch tuộc tối ưu bao gồm: Sơ chế tẩy nhớt bằng muối $\rightarrow$ Muối ướt 15-20\% $\rightarrow$ Sấy đối lưu $45-50^\circ C$ đến ẩm 18\% $\rightarrow$ Xông khói lạnh $30^\circ C$ trong 2-3 giờ $\rightarrow$ Đóng gói chân không. Việc tuân thủ các thông số $a_w$ và nhiệt độ là chìa khóa để đạt chất lượng cảm quan vượt trội và an toàn vệ sinh thực phẩm.

\newpage
\renewcommand{\bibname}{Tài liệu tham khảo}
\begin{thebibliography}{99}
\addcontentsline{toc}{section}{Tài liệu tham khảo}

\bibitem{lipid_storage} PMC. (2024). \textit{The Effect of Low Temperature Storage on the Lipid Quality of Fish}. \url{https://pmc.ncbi.nlm.nih.gov/articles/PMC11011431/}

\bibitem{octopus_texture} ResearchGate. (2012). \textit{Impact of Physical and Chemical Pretreatments on Texture of Octopus}. \url{https://www.researchgate.net/publication/229726846_Impact_of_Physical_and_Chemical_Pretreatments_on_Texture_of_Octopus}

\bibitem{octopus_drying} ResearchGate. (2023). \textit{Thin-layer drying kinetics and quality assessment of octopus}. \url{https://www.researchgate.net/publication/404359995_Thin-layer_drying_kinetics_and_quality_assessment_of_octopus}

\bibitem{octopus_oregano} ResearchGate. (2009). \textit{Combined effect of vacuum-packaging and oregano essential oil on Mediterranean octopus}. \url{https://www.researchgate.net/publication/23950398_Combined_effect_of_vacuum-packaging_and_oregano_essential_oil}

\bibitem{squid_acids} Gokoglu, N. (2023). \textit{Physicochemical and textural properties of squid treated with organic acids}. \url{https://www.ojs.unimal.ac.id/index.php/acta-aquatica/article/view/8952}

\bibitem{octopus_methods} Okuda, Y. \textit{History of octopus processing and methods of processing}. \url{https://waves-vagues.dfo-mpo.gc.ca/Library/18907.pdf}

\bibitem{octopus_patent} Google Patents. \textit{Octopus processing technology standard}. \url{https://patents.google.com/patent/CN104544326A/en}

\bibitem{wet_salting} ResearchGate. (2021). \textit{The effect of wet-salting preservation method on Dasyatis sp}. \url{https://www.researchgate.net/publication/353549559_The_effect_of_wet-salting_preservation_method}

\bibitem{aw_test} CTLA. \textit{Water Activity (aw) Test}. \url{https://ctlatesting.com/blogs/articles/water-activity-aw-test}

\bibitem{case_hardening} VTT. \textit{Analysis of Case Hardening}. \url{https://publications.vtt.fi/pdf/jurelinkit/RTE_Ranta-Maunus7.pdf}

\bibitem{vacuum_drying} PMC. (2019). \textit{Hot-Air Drying and Vacuum Microwave Drying of Cassia alata}. \url{https://pmc.ncbi.nlm.nih.gov/articles/PMC6515325/}

\bibitem{octopus_kinetics} Biotech Asia. (2021). \textit{Study of Salted Octopus Drying Kinetics}. \url{https://www.biotech-asia.org/vol18no3/study-of-salted-octopus-drying-kinetics/}

\bibitem{smoke_antioxidant} ResearchGate. \textit{Phenolic antioxidant in wood smoke}. \url{https://www.researchgate.net/publication/11761203_Phenolic_antioxidant_in_wood_smoke}

\bibitem{hun_khoi_ky_thuat} Hướng Nghiệp Á Âu. \textit{Kỹ Thuật Hun Khói Trên Nguyên Liệu Thịt, Cá}. \url{https://www.huongnghiepaau.com/ky-thuat-hun-khoi}

\bibitem{lychee_pulp} PMC. (2024). \textit{Phenolic Changes in Dried Lychee Pulp}. \url{https://pmc.ncbi.nlm.nih.gov/articles/PMC12896510/}

\bibitem{go_soi_xong_khoi} Vỏ xúc xích Collagen. \textit{Gỗ sồi xông khói nhập khẩu từ Đức}. \url{https://voxucxichcollagen.com.vn/go-soi-xong-khoi/}

\bibitem{tcvn_5649} HTVLAB. \textit{TCVN 5649:2006 về thủy sản khô - yêu cầu vệ sinh}. \url{https://htvlab.com/thong-tin-huu-ich/tieu-chuan/thuy-san-kho-yeu-cau-ve-sinh-tieu-chuan-viet-nam-tcvn-5649-2006.html}

\end{thebibliography}

\end{document}
